% GENERATED FILE. Edit canonical lessons, then run npm run generate:books.
% canonical-source-hash: fnv1a64:689d723e12d1eda3
% canonical-lessons: RU-C23-dva, RU-C23-tri, RU-C23-chetyre, RU-C23-pyat, RU-C23-practice

\chapter{One to Five, and an Answer to Skolko}
\label{ch:ru-count-five}
\hlchaptermodality{23}

\begin{chapteropening}
\textbf{By the end of this chapter:} \emph{I can say, hear, read and write the Russian numbers one to five, answer skolko with each of them, and name the English cousin of each.}

Five numbers, five Indo-European roots and five everyday English cousins -- the first answers this book can give to a question it has been able to ask since the last chapter.
\end{chapteropening}

\section[dva]{\ru{два} — two, and the question this book could already ask}

\label{lesson:RU-C23-dva}



\begin{quote}
Ask \textbf{\ru{сколько}} — \emph{how many}. Then say \textbf{\ru{один}} and the root it goes back to.
\end{quote}

\subsection*{You'll want to know first — one word}
\begin{quote}
\textbf{\ru{два}} — \emph{dva} — \textbf{two}
\end{quote}

Twenty-two chapters in, this book can ask \textbf{\ru{сколько}?} and could not understand any answer to it. \textbf{\ru{один}} is in the corpus, but it arrived as half of the \textbf{\ru{один} … \ru{другой}} pattern — \emph{one … the other} — which is a joining word wearing a number's clothes. \textbf{\ru{два}} is the first number this track teaches as a number.

\begin{cousinweb}
\textbf{\ru{два}} continues Indo-European *\emph{dwóh\textsubscript{1}}, \textquotedblleft{}two.\textquotedblright{} English \textbf{two}, Latin \textbf{duo}, Greek \textbf{dýo} and Sanskrit \textbf{dvá} are the same word, and Russian's form is among the least worn of them.

Say \emph{dva} and \emph{two} one after the other. That is not something to learn; it is something to recognise, and it will be true of nearly every number in this slice.
\end{cousinweb}

\subsection*{Your turn}
\begin{itemize}
\raggedright
  \item \emph{Say it:} \textbf{\ru{два}} — two
  \item \emph{Connect:} \textbf{\ru{два}} $\leftarrow$ *\emph{dwóh\textsubscript{1}} $\to$ English \textbf{two}, Latin \textbf{duo}
  \item \emph{Run through:} \textbf{\ru{один}, \ru{два}}
\end{itemize}

\begin{tcolorbox}[breakable,colback=teal!4,colframe=teal!35!black,title={Before you move on}]
Sources: \href{https://en.wiktionary.org/wiki/\%D0\%B4\%D0\%B2\%D0\%B0}{Wiktionary: \ru{два}}; \href{https://www.omniglot.com/language/numbers/russian.htm}{Omniglot, Numbers in Russian}.
\end{tcolorbox}
\section[tri]{\ru{три} — three}

\label{lesson:RU-C23-tri}



\begin{quote}
Say \textbf{\ru{два}} and its English cousin.
\end{quote}

\subsection*{You'll want to know first — one word}
\begin{quote}
\textbf{\ru{три}} — \emph{tri} — \textbf{three}
\end{quote}

Three letters, and every one of them is a Cyrillic shape this book taught long ago. Say it and read it at the same time: \textbf{\ru{т}-\ru{р}-\ru{и}}.

\begin{cousinweb}
\textbf{\ru{три}} continues Indo-European *\emph{tréyes}, \textquotedblleft{}three.\textquotedblright{} English \textbf{three}, Latin \textbf{trēs}, Greek \textbf{treîs} and Sanskrit \textbf{tri} are the same word.

This is the plainest of the whole set. Nothing has shifted: the \emph{t}, the \emph{r} and the vowel are all where they were, which is why \textbf{\ru{три}} and \textbf{three} sound like each other and \textbf{\ru{сто}} and \textbf{hundred}, later in this slice, will not.
\end{cousinweb}

\subsection*{Your turn}
\begin{itemize}
\raggedright
  \item \emph{Say it:} \textbf{\ru{три}} — three
  \item \emph{Connect:} \textbf{\ru{три}} $\leftarrow$ *\emph{tréyes} $\to$ English \textbf{three}, Latin \textbf{trēs}
  \item \emph{Run through:} \textbf{\ru{один}, \ru{два}, \ru{три}}
\end{itemize}

\begin{tcolorbox}[breakable,colback=teal!4,colframe=teal!35!black,title={Before you move on}]
Sources: \href{https://en.wiktionary.org/wiki/\%D1\%82\%D1\%80\%D0\%B8}{Wiktionary: \ru{три}}; \href{https://www.omniglot.com/language/numbers/russian.htm}{Omniglot, Numbers in Russian}.
\end{tcolorbox}
\section[chetýre]{\ru{четыре} — four, and one consonant's three fates}

\label{lesson:RU-C23-chetyre}



\begin{quote}
Say \textbf{\ru{три}} and its Latin cousin, then \textbf{\ru{два}}.
\end{quote}

\subsection*{You'll want to know first — one word}
\begin{quote}
\textbf{\ru{четыре}} — \emph{chetýre} — \textbf{four}
\end{quote}

Three syllables, stressed on the second, and the vowel in that stressed syllable is \textbf{\ru{ы}} — the sound this book warned has no English equivalent. Hold it there and the word falls into place.

\begin{cousinweb}
\textbf{\ru{четыре}} continues Indo-European *\emph{k\textsuperscript{w}etwóres}, \textquotedblleft{}four.\textquotedblright{}

The opening *\emph{k\textsuperscript{w}} is the interesting part, because the branches split three ways over it. Latin kept it whole: \textbf{quattuor}. English lost the lip-rounding and softened the \emph{k} to \emph{f}: \textbf{four}. Slavic pushed it forward to a \textbf{ch}, which is what you are saying.
\end{cousinweb}

\subsection*{Your turn}
\begin{itemize}
\raggedright
  \item \emph{Say it:} \textbf{\ru{четыре}} — four, stress on the \textbf{\ru{ы}}
  \item \emph{Connect:} \textbf{\ru{четыре}} $\leftarrow$ *\emph{k\textsuperscript{w}etwóres} $\to$ Latin \textbf{quattuor}, English \textbf{four}
  \item \emph{Run through:} \textbf{\ru{два}, \ru{три}, \ru{четыре}}
\end{itemize}

\begin{tcolorbox}[breakable,colback=teal!4,colframe=teal!35!black,title={Before you move on}]
Sources: \href{https://en.wiktionary.org/wiki/\%D1\%87\%D0\%B5\%D1\%82\%D1\%8B\%D1\%80\%D0\%B5}{Wiktionary: \ru{четыре}}; \href{https://www.omniglot.com/language/numbers/russian.htm}{Omniglot, Numbers in Russian}.
\end{tcolorbox}
\section[pyat]{\ru{пять} — five, and a cousin English kept whole}

\label{lesson:RU-C23-pyat}



\begin{quote}
Say \textbf{\ru{четыре}} and what three branches did with its first sound.
\end{quote}

\subsection*{You'll want to know first — one word}
\begin{quote}
\textbf{\ru{пять}} — \emph{pyat} — \textbf{five}
\end{quote}

One syllable, ending in the soft sign that this book has been putting on the ends of words since its first chapters. It softens the \textbf{\ru{т}} and is not a sound of its own.

\begin{cousinweb}
\textbf{\ru{пять}} continues Indo-European *\emph{pénk\textsuperscript{w}e}, \textquotedblleft{}five.\textquotedblright{} Greek \textbf{pénte}, Latin \textbf{quīnque} and English \textbf{five} are the same word.

The Greek form is the one an English speaker already owns: \textbf{pentagon}, \textbf{pentathlon}, \textbf{Pentecost} — all of them carry \emph{five} in a shape much closer to \textbf{\ru{пять}} than English \textbf{five} is. The number was in the reader's vocabulary before the lesson was.
\end{cousinweb}

\subsection*{Your turn}
\begin{itemize}
\raggedright
  \item \emph{Say it:} \textbf{\ru{пять}} — five
  \item \emph{Connect:} \textbf{\ru{пять}} $\leftarrow$ *\emph{pénk\textsuperscript{w}e} $\to$ Greek \textbf{pénte}, English \textbf{five}
  \item \emph{Run through:} \textbf{\ru{один}, \ru{два}, \ru{три}, \ru{четыре}, \ru{пять}} — the whole first five
\end{itemize}

\begin{tcolorbox}[breakable,colback=teal!4,colframe=teal!35!black,title={Before you move on}]
Sources: \href{https://en.wiktionary.org/wiki/\%D0\%BF\%D1\%8F\%D1\%82\%D1\%8C}{Wiktionary: \ru{пять}}; \href{https://www.omniglot.com/language/numbers/russian.htm}{Omniglot, Numbers in Russian}.
\end{tcolorbox}
\section[Practice]{Payoff — one to five, and an answer to \ru{сколько}}

\label{lesson:RU-C23-practice}



\begin{quote}
Ask \textbf{\ru{сколько}}, say \textbf{\ru{пять}}, then say what \textbf{\ru{один}} goes back to.
\end{quote}

\subsection*{Your turn}
\begin{enumerate}
\raggedright
  \item Hear each of the five named singly and out of order, and say what it is.
  \item Count backwards from \textbf{\ru{пять}} to \textbf{\ru{один}}.
  \item Read the five printed shuffled — \textbf{\ru{пять}, \ru{один}, \ru{четыре}, \ru{три}, \ru{два}} — and say each.
  \item Write all five from dictation. No new Cyrillic letter appears in any of them.
  \item Hear \textbf{\ru{сколько}?} and answer it with each of the five in turn.
\end{enumerate}

Then say the English cousin of each: \textbf{one, two, three, four, five}. All five have one. That is not luck — Russian's numerals are among the least worn in the family, and this book's reader already owned most of them.

\begin{tcolorbox}[breakable,colback=teal!4,colframe=teal!35!black,title={Before you move on}]
Source: \href{https://www.omniglot.com/language/numbers/russian.htm}{Omniglot, Numbers in Russian}.
\end{tcolorbox}
